
\documentclass[UTF8,11pt,a4paper]{ctexart}
\usepackage[a4paper,left=23mm,right=23mm,top=24mm,bottom=24mm]{geometry}
\usepackage{fontspec}
\setmainfont{DejaVu Sans}
\setCJKmainfont{Noto Serif CJK SC}
\usepackage{amsmath,amssymb}
\usepackage{xcolor}
\usepackage[colorlinks=true,linkcolor=blue!45!black,unicode=true]{hyperref}
\usepackage{enumitem}
\usepackage{needspace}
\usepackage{fancyhdr}
\setlength{\parindent}{2em}
\setlength{\parskip}{0.25em}
\setlist[enumerate]{leftmargin=2.6em,itemsep=0.4em,topsep=0.5em}
\pagestyle{fancy}
\fancyhf{}
\fancyhead[L]{\small 三千弱水 · 大学化学高阶综合题库}
\fancyhead[R]{\small 待专家审校稿}
\fancyfoot[C]{\thepage}
\setlength{\headheight}{15pt}
\sloppy

\begin{document}
\begin{center}
{\LARGE\bfseries 大学化学高阶综合题 · 第3辑 (12题)}\\[0.8em]
{\large 题目、分步解析与数值核验}\\[0.5em]
{\small \textcolor{red!65!black}{审校稿：未经化学专家核准；禁止直接视为正式试卷或出版物}}
\end{center}
\vspace{1em}
\noindent 本册从仓库结构化 JSONL 自动生成。增补和修改应只编辑源题目数据，并重新执行统一编译。
\par 机械方程残差检查只能说明指定模型的内部数值一致，不保证模型现实可行性。第六辑另须完成跨辑模型去重审查。
\tableofcontents
\clearpage
\section*{第3辑 · 题目及解答}
\addcontentsline{toc}{section}{第3辑}
\subsection*{无机化学}
\addcontentsline{toc}{subsection}{第3辑 / 无机化学}
\Needspace{7\baselineskip}
\subsubsection*{UC-ADV3-I04 · 二级络合与配体质子化的闭合物料平衡}
\noindent\textbf{题目}\par
假想二价金属M的总浓度为0.01200 mol/L，单酸配体HL的总分析浓度为0.03000 mol/L；固定pH=8.000且pKa(HL)=8.000。只允许M2+、ML+、ML2及游离HL、L−存在，逐级形成常数K1=250 L/mol、K2=80 L/mol，所有常数均采用理想浓度模型，体积恒定且缓冲液及电中性反离子由外部提供。求游离L−、HL、M2+、ML+、ML2的浓度，金属的双络合分数及被络合配体的比例。不得把配体总浓度直接当作游离L−。最后证明物理解唯一，并用配体与金属双重守恒核验。\par
\medskip\noindent\textbf{分步解答}
\begin{enumerate}
\item 酸碱关系Ka=[H+][L−]/[HL]；当pH=pKa，游离HL与L−浓度相等。令l=[L−]，则[HL]=l。
\item 令m=[M2+]。逐级络合给[ML+]=K1ml、[ML2]=K1K2ml²。
\item 金属守恒CM=m(1+K1l+K1K2l²)，故m(l)=0.01200/(1+250l+20000l²)。
\item 配体守恒CL=[HL]+[L−]+[ML+]+2[ML2]，即F(l)=2l+0.01200(250l+40000l²)/(1+250l+20000l²)−0.03000=0。
\item F(l)在l≥0严格递增：每个金属平均结合配体数作为自由配体活度的函数严格递增，且2l也严格递增；F(0)<0、F(0.03000)>0，故唯一正根l=0.008438617122 mol/L。
\item 回代[HL]=0.008438617122、[M2+]=0.002646751649、[ML+]=0.005583730946、[ML2]=0.003769517405 mol/L。
\item 双络合金属分数=[ML2]/CM=0.3141264504；被络合配体比例=([ML+]+2[ML2])/CL=0.4374255252。
\item 检验金属守恒m+[ML+]+[ML2]=0.012 mol/L，配体守恒2l+[ML+]+2[ML2]=0.03 mol/L；若误取l=CL，两条守恒不可能同时成立。
\end{enumerate}
\noindent\textbf{答案摘要}\par
答案1 = 0.00843861712 mol/L；答案2 = 0.00843861712 mol/L；答案3 = 0.00264675165 mol/L；答案4 = 0.00558373095 mol/L；答案5 = 0.00376951741 mol/L；答案6 = 0.31412645 fraction；答案7 = 0.437425525 fraction\par
\noindent\textbf{模型前提：}
由外部缓冲体系固定pH；忽略其它络合物、活度修正和缓冲液中可能的竞争配位；不宣称给出具体真实金属的实验形成常数。\par
\medskip\hrule\medskip
\Needspace{7\baselineskip}
\subsubsection*{UC-ADV3-I05 · 氧空位供体与受主掺杂的电子—空穴电中性反演}
\noindent\textbf{题目}\par
一假想n型氧化物在300 K下氧位密度为4.00×10\textasciicircum{}22 cm\textasciicircum{}-3，带+2有效电荷的氧空位浓度[V\_O\textasciicircum{}\{••\}]=1.00×10\textasciicircum{}17 cm\textasciicircum{}-3，完全电离的固定受主掺杂浓度[A′]=5.00×10\textasciicircum{}16 cm\textasciicircum{}-3。电子与空穴分别带−1和+1有效电荷，且处于非简并热平衡，np=ni²、ni=1.00×10\textasciicircum{}10 cm\textasciicircum{}-3。电子迁移率20.0 cm²/(V s)，空穴迁移率5.00 cm²/(V s)，基本电荷1.602176634×10\textasciicircum{}-19 C。求n、p、氧空位位点分数、总电子电导率与空穴贡献比例；证明所取代数根是唯一物理解，并解释为何不能仅写n=[V\_O]。\par
\medskip\noindent\textbf{分步解答}
\begin{enumerate}
\item Kröger–Vink有效电荷守恒：p+2[V\_O\textasciicircum{}\{••\}]=n+[A′]，即n−p=2×10\textasciicircum{}17−5×10\textasciicircum{}16=150000000000000000 cm\textasciicircum{}-3。
\item 与本征热平衡np=ni²联立，消去p得到n²−Δn−ni²=0，其中Δ=1.50×10\textasciicircum{}17 cm\textasciicircum{}-3。
\item 二次方程正根n=(Δ+√(Δ²+4ni²))/2=150000000000000000 cm\textasciicircum{}-3；负根为负电子浓度，必须舍去。
\item 为避免极小空穴浓度出现数值抵消，使用p=ni²/n=666.6666667 cm\textasciicircum{}-3，而不是两个大数直接相减。
\item 氧空位位点分数fV=[V\_O]/N\_O=0.0000025，即2.5 ppm。
\item 电导率σ=e(μn n+μp p)=0.4806529902 S/cm；空穴贡献比例μp p/(μn n+μp p)=1.111111111e-15。
\item 核验n−p=150000000000000000 cm\textasciicircum{}-3，np=100000000000000000000 cm\textasciicircum{}-6；因此约两电子/空位但还需扣除受主补偿，不能机械设n=[V\_O]。
\end{enumerate}
\noindent\textbf{答案摘要}\par
答案1 = 1.5e+17 cm\textasciicircum{}-3；答案2 = 666.666667 cm\textasciicircum{}-3；答案3 = 2.5e-06 fraction；答案4 = 0.48065299 S/cm；答案5 = 1.11111111e-15 fraction\par
\noindent\textbf{模型前提：}
所有缺陷浓度固定且完全电离，载流子非简并、迁移率给定，无缺陷关联或电极化学反应；这是缺陷化学的假想简化模型。\par
\medskip\hrule\medskip
\Needspace{7\baselineskip}
\subsubsection*{UC-ADV3-I06 · 两性金属氢氧化物溶解度的非单调性与全局极小值}
\noindent\textbf{题目}\par
过量固态M(OH)2与强pH缓冲溶液平衡，假设溶解只形成M2+及[M(OH)4]2−，不允许其它溶解物种。Ksp=[M2+][OH−]\textasciicircum{}2=1.00×10\textasciicircum{}-16（浓度型），总形成常数β4=[M(OH)4]2−/([M2+][OH−]\textasciicircum{}4)=1.00×10\textasciicircum{}10 L\textasciicircum{}4 mol\textasciicircum{}-4。在25°C取pKw=14.000，理想活度近似，背景反离子电荷由缓冲体系承担。(1)建立总金属溶解度s(pH)。(2)严格求使s最小的OH−浓度及pH，证明是全局最小。(3)计算pH=11.000、12.000的溶解度并区分低pH游离离子与高pH羟基络合支配区。(4)做沉淀平衡与络合平衡的逆向核对。\par
\medskip\noindent\textbf{分步解答}
\begin{enumerate}
\item 令u=[OH−]>0，沉淀平衡给[M2+]=Ksp/u²，形成常数给[M(OH)4]2−=β4 Ksp u²。
\item 总溶解度s(u)=Ksp u\textasciicircum{}(−2)+β4 Ksp u²；前项随u增大而下降、后项上升。
\item 一阶导数s′(u)=−2Ksp u\textasciicircum{}(−3)+2β4 Ksp u；令s′=0，得u⁴=1/β4。
\item 正根u*=β4\textasciicircum{}(−1/4)=0.00316227766 mol/L；s″=6Ksp/u⁴+2β4Ksp>0，且u→0或∞时s→∞，故是唯一全局最小。
\item pOH*=−log10(u*)=2.5，pH*=14−pOH*=11.5；smin=2Ksp√β4=2e-11 mol/L。
\item pH=11时u=10\textasciicircum{}-3 M，s11=1.01e-10 mol/L；pH=12时u=10\textasciicircum{}-2 M，s12=1.01e-10 mol/L。
\item 在最小点游离M2+与四羟基络离子浓度均为1e-11 mol/L；将两个物种代入Ksp与β4公式分别恢复1.00×10\textasciicircum{}-16和1.00×10\textasciicircum{}10。
\end{enumerate}
\noindent\textbf{答案摘要}\par
答案1 = 0.00316227766 mol/L；答案2 = 11.5 pH；答案3 = 2e-11 mol/L；答案4 = 1.01e-10 mol/L；答案5 = 1.01e-10 mol/L\par
\noindent\textbf{模型前提：}
题设仅适用于具有该指定形成常数的假想两性氢氧化物；极端pH可能出现其他多核物种、活度与缓冲副反应，本模型已明确排除。\par
\medskip\hrule\medskip
\subsection*{分析化学}
\addcontentsline{toc}{subsection}{第3辑 / 分析化学}
\Needspace{7\baselineskip}
\subsubsection*{UC-ADV3-A04 · 二次基质响应中的三点标准加入与歧义根排除}
\noindent\textbf{题目}\par
一份未知分析物的原液在定容后放入相同基质的三个测量池，分别加入浓度增量s=0、0.000800、0.001600 mol/L，三池终体积相同。校准空白信号b=0.015000已知；真实测量规律A(s)=b+k(c+s)+q(c+s)²，其中未知c>0、k>0、q>0（后两者分别有适用的量纲），且分析物先验浓度范围0<c<0.005 mol/L。测得信号依次A0=0.054845，A1=0.080205，A2=0.106205（按题示有效位数视为精确）。(1)由三点差分求二次响应系数、表观一阶系数及截距。(2)将原浓度c作为二次方程的根，求两代数根并依据物理约束舍去伪根。(3)复原k、q并验证三点信号，讨论为什么仅两点不足以分离三个未知量。\par
\medskip\noindent\textbf{分步解答}
\begin{enumerate}
\item 将响应写成观测加标增量s的多项式A(s)=A0+B s+q s²，其中B=k+2qc，A0−b=kc+qc²。
\item 等间隔二阶差分Δ²A=A2−2A1+A0=2q(Δs)²，得到q=500 L²/mol²。
\item 第一差分(A1−A0)/Δs=B+qΔs，得B=31.3 L/mol。
\item 记C=A0−b=0.039845，由k=B−2qc代回C=kc+qc²得到q c²−B c+C=0。
\item 判别式D=B²−4qC=900>0，两个形式根分别c小=0.0013 mol/L与c大=0.0613 mol/L。
\item 按先验c<0.005 mol/L且k>0，只允许c=0.0013；另一根超出适用浓度区并导致k<0。
\item 恢复k=B−2qc=30 L/mol，q=500 L²/mol²；逐一代入测量方程得A(0)=0.054845，A(Δs)=0.080205，A(2Δs)=0.106205。
\item 若只测两个加标水平，仅给两条独立方程而c、k、q有三个未知量，因此在缺少先验或其他信息时不可辨识。
\end{enumerate}
\noindent\textbf{答案摘要}\par
答案1 = 0.0013 mol/L；答案2 = 30 L/mol；答案3 = 500 L²/mol²；答案4 = 0.0613 mol/L\par
\noindent\textbf{模型前提：}
三次测量使用相同基质、完全相同最终体积；响应在目标浓度区确实为二次多项式，所给数据视为理想精确值，不据此声称真实仪器无限精度。\par
\medskip\hrule\medskip
\Needspace{7\baselineskip}
\subsubsection*{UC-ADV3-A05 · 同位素稀释与化学回收率的双信号联合反演}
\noindent\textbf{题目}\par
取100.0 mL未知氯化物水样，原有氯的37Cl同位素原子分数x天然=0.2422。向其加入总氯量ns=0.000800 mol、37Cl分数x标=0.9800的富集同位素内标，充分混合后进行化学分离。分离步骤对两种同位素具有相同且未知的回收率r。质谱已校正质量歧视，测得分离后37Cl原子分数x=0.453；总氯信号S=4760计数，独立绝对信号斜率α=2.000×10\textasciicircum{}6 计数/mol已标定。求原始样品氯分析浓度、共同回收率与测量的同位素物料守恒。最后计算在x附近样品物质的量对x的导数dn/dx（mol/单位分数），解释同位素比值为何能与总回收率解耦。\par
\medskip\noindent\textbf{分步解答}
\begin{enumerate}
\item 设原样氯物质的量为n，混合前样品和内标氯量分别为n、ns；回收后两者同乘r。
\item 同位素原子分数x=[r x天然 n+r x标 ns]/[r(n+ns)]=(x天然 n+x标 ns)/(n+ns)，共同回收率r在比例中消去。
\item 整理得n=ns(x标−x)/(x−x天然)=0.002 mol，原水样分析浓度c=n/0.1000 L=0.02 mol/L。
\item 绝对信号S=α r(n+ns)；用已由同位素比值求出的n反算回收率r=S/[α(n+ns)]=0.85。
\item 回收后样品氯量rn=0.0017 mol，内标氯量rns=0.00068 mol；实际总回收量=0.00238 mol。
\item 把两种已回收物质的37Cl量相加再除以总回收量，恢复观测x=0.453，且以α乘总回收量恢复S=4760计数。
\item 由n=ns(x标−x)/(x−x天然)求导，dn/dx=−ns(x标−x天然)/(x−x天然)²=-0.01328273245 mol/单位同位素分数。
\item 当x接近天然丰度时，分母变小导致反演对测量误差敏感；如果两种同位素的化学回收率不同，则共同r消去的关键步骤不成立。
\end{enumerate}
\noindent\textbf{答案摘要}\par
答案1 = 0.002 mol；答案2 = 0.02 mol/L；答案3 = 0.85 fraction；答案4 = -0.0132827324 mol per isotope-fraction\par
\noindent\textbf{模型前提：}
假设无同位素交换/分馏、无其它氯来源、内标与原样完全混合；绝对响应斜率已独立标定，计数视作无随机波动的模型输入。\par
\medskip\hrule\medskip
\Needspace{7\baselineskip}
\subsubsection*{UC-ADV3-A06 · 双弱酸物种分布与跨pH光谱的耦合解析}
\noindent\textbf{题目}\par
未知水样仅含单酸HA、HB两种分析物，总浓度分别记CA、CB，在两个独立的强缓冲测量池中精确控制pH为6.000和8.000；每池通过定容确保分析物总浓度均与原液相同。HA的Ka=1.00×10\textasciicircum{}-5，HB的Ka=1.00×10\textasciicircum{}-7，仅其阴离子A−、B−显色，未解离中性酸不吸光。1.00 cm池、两波长阴离子摩尔吸光系数矩阵为[[12000,2000],[3000,10000]] L mol\textasciicircum{}-1 cm\textasciicircum{}-1。已测λ1在pH6的吸光度0.3309090909和λ2在pH8的吸光度0.2717282717，另有λ2在pH6的独立检验值0.1。(1)以分布系数推导2×2联立方程并求CA、CB。(2)预报第三个吸光度并核验。(3)解释为何不能把两种弱酸都按完全解离处理，说明两个波长/两个pH提供独立信息的矩阵可逆条件。\par
\medskip\noindent\textbf{分步解答}
\begin{enumerate}
\item 对单酸HA ⇌H++A−，在固定pH下阴离子分数αA(pH)=KaA/[KaA+10\textasciicircum{}(−pH)]；HB同理。
\item 在pH6时αA=0.9090909091、αB=0.09090909091；在pH8时αA=0.999000999、αB=0.9090909091。
\item Beer–Lambert关系为Aλ1,pH6=(12000αA6)CA+(2000αB6)CB=0.3309090909。
\item 另一个独立关系Aλ2,pH8=(3000αA8)CA+(10000αB8)CB=0.2717282717。
\item 系数矩阵为[[10909.09091,181.8181818],[2997.002997,9090.909091]]，行列式D=98628644.08≠0，故两个总浓度唯一确定。
\item 按克拉默法则CA=(m22 A1−m12 A2)/D=0.00003 mol/L，CB=(m11 A2−m21 A1)/D=0.00002 mol/L。
\item 不参与求解的第三信号理论值Aλ2,pH6=3000αA6CA+10000αB6CB=0.1，与独立测得0.1一致。
\item 若错误假设两酸在pH6完全解离，会高估较弱的HB离子分数并产生系统误差；可辨识性取决于D≠0而非仅有两条读数。
\end{enumerate}
\noindent\textbf{答案摘要}\par
答案1 = 3e-05 mol/L；答案2 = 2e-05 mol/L；答案3 = 0.1 absorbance；答案4 = 98628644.1 L²/mol²\par
\noindent\textbf{模型前提：}
酸碱平衡受外部强缓冲维持，测量时不发生化学降解或体积变化；吸光定律、两阴离子光谱系数已标定，忽略高浓度活度与交互作用。\par
\medskip\hrule\medskip
\subsection*{有机化学}
\addcontentsline{toc}{subsection}{第3辑 / 有机化学}
\Needspace{7\baselineskip}
\subsubsection*{UC-ADV3-O04 · 自催化异构化Logistic方程、最优速率与目标转化时间}
\noindent\textbf{题目}\par
恒容恒温假想同分异构化A+B→2B遵循速率律r=k[A][B]，k=5.000 L mol\textasciicircum{}-1 h\textasciicircum{}-1，初始[A]=0.08000、[B]=0.002000 mol/L。仅有该反应，A、B分子式及摩尔质量相同。求B(t)的解析式；B浓度升至0.06000 mol/L所需时间及A转化率；证明反应速率何时唯一达到峰值并求峰值和当时浓度；用总物料守恒与逆代入核对。题目严禁将B近似看成恒定催化剂。\par
\medskip\noindent\textbf{分步解答}
\begin{enumerate}
\item 化学计量为消耗1个A净生成1个B，故C=[A]+[B]=0.082 mol/L恒定。
\item 令b=[B]，则db/dt=k(C−b)b；由于b同时是催化剂与产物，得到非线性Logistic微分方程。
\item 由分离变量与初始条件b(0)=0.002000，b(t)=C/[1+(0.08000/0.002000)e\textasciicircum{}(−kCt)]。
\item 反解时间t(b)=ln\{[b/(C−b)]/[B0/A0]\}/(kC)，得到b=0.06000时t=11.44434528 h。
\item 由b−B0为A转化量，X\_A=(0.06000−0.002000)/0.08000=0.725；此时[A]=0.022 mol/L。
\item 将速率看成b的函数r=kb(C−b)，开口向下的二次多项式于b*=C/2=0.041 mol/L唯一达到最大。
\item 该速率最大时刻t*=8.997266961 h，rmax=kC²/4=0.008405 mol L\textasciicircum{}-1 h\textasciicircum{}-1。
\item 逆代入b(t目标)=0.06、b(t*)=0.041；两处均保持[A]+[B]=C，且t>0时b单调上升。
\end{enumerate}
\noindent\textbf{答案摘要}\par
答案1 = 11.4443453 h；答案2 = 0.725 fraction；答案3 = 8.99726696 h；答案4 = 0.041 mol/L；答案5 = 0.008405 mol/(L h)\par
\noindent\textbf{模型前提：}
理想均相不可逆自催化异构化模型，恒定k、无副反应、无体积变化及催化失活。\par
\medskip\hrule\medskip
\Needspace{7\baselineskip}
\subsubsection*{UC-ADV3-O05 · 平行产物Eyring选择性反演与有限转化时间}
\noindent\textbf{题目}\par
假想有机不可逆平行一级反应A→P1、A→P2，两路共享相同的Eyring前因子kBT/h。第一路径已知活化焓95.00 kJ/mol、活化熵−30.00 J/(mol K)，300 K与340 K的产物初始速率比k1/k2依次为2.229585515和1.761186189，均视作理想精确读数。取R=8.314462618 J/(mol K)，kB=1.380649×10\textasciicircum{}-23 J/K，h=6.62607015×10\textasciicircum{}-34 J s，ki=(kBT/h)exp(ΔSi‡/R−ΔHi‡/(RT))。反演第二路径活化焓与熵；在320 K计算两速率常数、产物选择性、初始纯A总转化80\%所需时间及两产物的实际收率；检查Eyring数据与总体物料守恒。\par
\medskip\noindent\textbf{分步解答}
\begin{enumerate}
\item 定义r=k1/k2，Eyring相除得lnr=ΔΔS‡/R−ΔΔH‡/(RT)，其中差量是路径1减路径2。
\item 由两温度比值消去熵差：ΔΔH‡=−R ln(r340/r300)/(1/340−1/300)=-5 kJ/mol。
\item ΔΔS‡=Rlnr300+ΔΔH‡/300=-10 J/(mol K)。
\item 据此第二路径ΔH2‡=95.00−ΔΔH‡=100 kJ/mol，ΔS2‡=−30.00−ΔΔS‡=-20 J/(mol K)。
\item 直接代入320 K得k1=0.00005624872133 s\textasciicircum{}-1、k2=0.00002859555573 s\textasciicircum{}-1，产物分支比k1/k2=1.967044175。
\item 两条一级路径共有初始A，A(t)/A0=exp[−(k1+k2)t]，80\%转化对应t80=ln5/(k1+k2)=18969.31612 s。
\item 两产物相对初始A的收率分别为Y1=0.8k1/(k1+k2)=0.5303713889和Y2=0.2696286111，未反应A为0.2000。
\item 回代温度300与340 K的Eyring速率比恢复原观测；Y1+Y2+0.2000=1，严格满足物料守恒。
\end{enumerate}
\noindent\textbf{答案摘要}\par
答案1 = 100000 J/mol；答案2 = -20 J/(mol K)；答案3 = 5.62487213e-05 s\textasciicircum{}-1；答案4 = 2.85955557e-05 s\textasciicircum{}-1；答案5 = 18969.3161 s；答案6 = 0.530371389 fraction；答案7 = 0.269628611 fraction\par
\noindent\textbf{模型前提：}
两路径严格一级、不可逆、常温区间内ΔH‡及ΔS‡为常量，传输系数及前因子相同，无产物副反应。\par
\medskip\hrule\medskip
\Needspace{7\baselineskip}
\subsubsection*{UC-ADV3-O06 · 动力学拆分中剩余ee与总转化率的逆问题}
\noindent\textbf{题目}\par
某假想手性原料初始为严格外消旋R/S各0.5000 mol。在恒温封闭体系中，两对映体以不同的一级速率常数不可逆生成各自的PR/PS，无消旋、交叉反应或其它产物。10.000 h后测得底物总转化率X=0.4789937837，剩余底物ees=(nS−nR)/(nS+nR)=0.4218990053。反演两对映体一级速率常数、10h时产物的各自摩尔数与对映体过量eep=(nPR−nPS)/(nPR+nPS)；进一步求90\%总转化时的时间与eep。要求给出物理解唯一性说明、指数动力学回代与总物料守恒。\par
\medskip\noindent\textbf{分步解答}
\begin{enumerate}
\item 转化率定义X=1−(nR+nS)，故10h剩余底物总量为1−X=0.5210062163 mol。
\item 从ees=(nS−nR)/(nS+nR)，解得nR=(1−X)(1−ees)/2=0.150597106 mol、nS=(1−X)(1+ees)/2=0.3704091103 mol。
\item 独立一级模型nR=0.5exp(−kR t)、nS=0.5exp(−kS t)，得到kR=−ln(2nR)/10=0.12 h\textasciicircum{}-1、kS=0.03 h\textasciicircum{}-1。
\item 10h产物量nPR=0.5−nR=0.349402894 mol，nPS=0.5−nS=0.1295908897 mol。
\item 相应产物ee\_p=(nPR−nPS)/(nPR+nPS)=0.4589036682，可与剩余底物ee\_s正负方向作对照。
\item 达到90.00\%转化时剩余总量0.1000 mol，故必须满足exp(−kR t)+exp(−kS t)=0.2000。
\item 其左侧对t≥0严格递减，由初值2降至0，故时间根唯一；二分求得t90=53.90742857 h。
\item 此时剩余R=0.0007754720561、S=0.09922452794 mol，产物ee\_p=0.1093878399；反代两指数函数及nR+nS=0.1 mol确认90\%转化。
\end{enumerate}
\noindent\textbf{答案摘要}\par
答案1 = 0.12 h\textasciicircum{}-1；答案2 = 0.03 h\textasciicircum{}-1；答案3 = 0.458903668 fraction；答案4 = 53.9074286 h；答案5 = 0.10938784 fraction\par
\noindent\textbf{模型前提：}
底物严格外消旋投料，反应均为不可逆一级、没有消旋或产物相互转化，分析信号对两对映体有已校准响应因子。\par
\medskip\hrule\medskip
\subsection*{物理化学}
\addcontentsline{toc}{subsection}{第3辑 / 物理化学}
\Needspace{7\baselineskip}
\subsubsection*{UC-ADV3-P04 · 双后缀Margules液相的共沸点、汽液组成与混合稳定性}
\noindent\textbf{题目}\par
理想气相二元A/B在恒温下的纯组分饱和蒸气压分别PA*=80.00 kPa、PB*=40.00 kPa；液相服从二后缀Margules方程lnγA=1.500(1−xA)²，lnγB=1.500xA²。采用改进Raoult定律yiP=xiγiPi*，两组分总液相摩尔分数和为1，取正压、0<xA<1。(1)推导共沸条件，求是否存在内部共沸点及其组成、压力。(2)求xA=0.4000时泡点总压及初生蒸汽yA。(3)证明该Margules模型在全部组成范围的局部混合稳定性，并论证为何共沸点不是随意挑选的一个极值。\par
\medskip\noindent\textbf{分步解答}
\begin{enumerate}
\item 采用改进Raoult定律，泡点压力P(x)=xγA PA*+(1−x)γB PB*，蒸气组成yA=xγA PA*/P。
\item 内部共沸条件是蒸汽、液体组分相同，即yA=x；等价于γA PA*=γB PB*。
\item 对两后缀Margules式取对数：A[(1−x)²−x²]=A(1−2x)=ln(PB*/PA*)。
\item 因此xaz=[1+ln(PA*/PB*)/A]/2=0.7310490602，满足0<xaz<1，内部共沸存在且此式根唯一。
\item 将xaz带回：γA=1.11460704，γB=2.229214079，泡点共沸压力Paz=89.16856318 kPa，计算yaz=0.7310490602。
\item 在xA=0.4000时，泡点压力P=85.4221992 kPa；初生汽相yA=0.6428331289，相对于xA的富集差=0.2428331289。
\item 理想混合熵与过剩Gibbs能给g\_mix/(RT)=xlnx+(1−x)ln(1−x)+Ax(1−x)，曲率d²g/dx²=1/x+1/(1−x)−2A。
\item 因为1/x+1/(1−x)≥4>2A=3，故整个0<x<1范围局部稳定；共沸由汽液化学势相等与y=x决定，而非任意调参取压力极值。
\end{enumerate}
\noindent\textbf{答案摘要}\par
答案1 = 0.73104906 liquid mole fraction；答案2 = 89.1685632 kPa；答案3 = 85.4221992 kPa；答案4 = 0.642833129 vapor mole fraction；答案5 = 2.08604749 dimensionless curvature\par
\noindent\textbf{模型前提：}
气相理想、液相只使用二后缀Margules且处于单液相局部稳定区，无蒸气缔合、反应或压力导致的液相非理想压缩效应。\par
\medskip\hrule\medskip
\Needspace{7\baselineskip}
\subsubsection*{UC-ADV3-P05 · 含惰性气体的合成氨反应平衡、总压效应与进度唯一性}
\noindent\textbf{题目}\par
在固定温度下，理想气相平衡N2+3H2⇌2NH3，给定标准压力p°=1.000 bar、无量纲热力学平衡常数Kp=0.5000。初始取n(N2)=1.000 mol、n(H2)=3.000 mol、n(NH3)=0、惰性气体nI=2.000 mol，在总压恒定100.0 bar下达到平衡。(1)严格由分压和标准态定义写出反应商Qp(ξ)，包括反应引起的总摩尔数变化。(2)在物理区间0<ξ<1上求唯一平衡进度、氨摩尔数和摩尔分数。(3)若温度及投料不变但总压力改为50.0 bar，重新求平衡并解释压力效应。(4)通过Qp、非负组成、总物料与方程单调性核验；不得把Kp当作未经标准化的分压乘积。\par
\medskip\noindent\textbf{分步解答}
\begin{enumerate}
\item 设正向反应进度ξ，平衡摩尔数分别nN2=1−ξ，nH2=3−3ξ，nNH3=2ξ，nI=2，故n总=6−2ξ。
\item 令yi=ni/(6−2ξ)，理想气分压pi=yiP；标准化反应商Qp=(pNH3/p°)²/[(pN2/p°)(pH2/p°)³]。
\item 代入化学计量并化简：Qp(ξ,P)=4ξ²(6−2ξ)²/[27(1−ξ)\textasciicircum{}4]×(p°/P)²；各压力统一以bar为数值单位。
\item 固定P=100 bar，解Qp=0.5000，物理解ξ100=0.8585250533 mol，NH3=1.717050107 mol，总摩尔数=4.282949893 mol，氨摩尔分数=0.4009036177。
\item lnQp对ξ的导数为2/ξ−4/(6−2ξ)+4/(1−ξ)，在0<ξ<1为正；Qp由0连续增至∞，故任何正K均只有一个物理根。
\item 固定P=50 bar重解得到ξ50=0.8039372876 mol、yNH3,50=0.3660812066；可见降低压力降低平衡转化，因为反应气体总摩尔数减少2。
\item 对100bar回代Qp=0.5，50bar回代Qp=0.5；两次剩余H2、N2均非负，nN2+nH2+nNH3+nI=6−2ξ。
\item 不能把惰性气体排除在n总之外，它改变反应物与产物分压；也不能在两种不同压力下沿用同一个ξ。
\end{enumerate}
\noindent\textbf{答案摘要}\par
答案1 = 0.858525053 mol reaction extent；答案2 = 0.803937288 mol reaction extent；答案3 = 0.400903618 mole fraction；答案4 = 0.366081207 mole fraction\par
\noindent\textbf{模型前提：}
总压恒定但可变容的等温理想气体系，Kp不随压力变且使用p°=1bar，不考虑真实高压活度系数和反应速率限制。\par
\medskip\hrule\medskip
\Needspace{7\baselineskip}
\subsubsection*{UC-ADV3-P06 · 非谐振动基频与热带联合反演分子振动常数及温度}
\noindent\textbf{题目}\par
某假想双原子分子按Morse型非谐振动模型G(v)=ωe(v+1/2)−ωexe(v+1/2)²，G和波数单位均为cm\textasciicircum{}-1。测得v=0→1基频1950.0 cm\textasciicircum{}-1、v=1→2热带1900.0 cm\textasciicircum{}-1；已将谱线波数及仪器响应归一化后的积分强度比定义为I12/I01=2 N1/N0=0.08000，其中线强度比2来自规定的振子矩阵元平方(v+1)，并假设振动态处于热平衡。取hc/kB=1.438776877 cm K。(1)反求ωe与ωexe。(2)从热带/基频强度比求振动温度T。(3)预测0→2跃迁能量对应波数与2→3热带波数，说明为何仅靠该振子近似不能预测实际0→2跃迁强度。(4)用能级差和Boltzmann关系逆向检验。\par
\medskip\noindent\textbf{分步解答}
\begin{enumerate}
\item 相邻振动能级差为G(v+1)−G(v)=ωe−2ωexe(v+1)。
\item 因此ν01=ωe−2ωexe=1950 cm\textasciicircum{}-1，ν12=ωe−4ωexe=1900 cm\textasciicircum{}-1；两式相减得ωexe=(ν01−ν12)/2=25 cm\textasciicircum{}-1。
\item 回代得ωe=ν01+2ωexe=2000 cm\textasciicircum{}-1，属于正的非谐修正。
\item 已按题设规范化线强度I12/I01=2(N1/N0)=0.08000，因此N1/N0=0.04。
\item 热平衡Boltzmann分布N1/N0=exp[−(hc/kB)ν01/T]，求得T=(hc/kB)ν01/ln(2/0.08000)=871.6132783 K。
\item 跨两级的0→2跃迁能级差ν02=G(2)−G(0)=ν01+ν12=3850 cm\textasciicircum{}-1；而2→3相邻热带ν23=ωe−6ωexe=1850 cm\textasciicircum{}-1。
\item 将求得ωe、ωexe代回G(v)，得G1−G0=1950、G2−G1=1900 cm\textasciicircum{}-1；再代回Boltzmann式恢复0.08000强度比。
\item 谐振子电偶极选择定则通常只允许Δv=±1；0→2是否有实际强度取决于振动波函数及偶极矩函数的非谐项，本题能算跃迁能量却不能凭已给两条谱线确定其绝对强度。
\end{enumerate}
\noindent\textbf{答案摘要}\par
答案1 = 2000 cm\textasciicircum{}-1；答案2 = 25 cm\textasciicircum{}-1；答案3 = 871.613278 K；答案4 = 3850 cm\textasciicircum{}-1；答案5 = 1850 cm\textasciicircum{}-1\par
\noindent\textbf{模型前提：}
Morse型振动态、处于Boltzmann热平衡，归一化强度律与矩阵元平方线性规则是题定模型，不假定光谱实验通用的所有线形因子都自动相消。\par
\medskip\hrule\medskip
\end{document}
